\chapter*{ضمیمهٔ ۴: واقع‌گرایی و آرمان‌گرایی در سیاست}
\addcontentsline{toc}{chapter}{ضمیمهٔ ۴: واقع‌گرایی و آرمان‌گرایی}
\markboth{ضمیمهٔ ۴: واقع‌گرایی و آرمان‌گرایی}{ضمیمهٔ ۴: واقع‌گرایی و آرمان‌گرایی}

\begin{tcolorbox}[
    enhanced,
    colback=lightbg,
    colframe=darkgray,
    boxrule=2pt,
    arc=8pt,
    fonttitle=\bfseries\large,
    title={\rl{پرسش محوری}}
]
\begin{center}
\Large\textit{آیا سیاست باید بر اساس «آنچه هست» باشد یا «آنچه باید باشد»؟}
\end{center}
\end{tcolorbox}

\section*{مقدمه: دو نگاه به سیاست}

در تاریخ اندیشهٔ سیاسی، دو رویکرد بنیادین وجود داشته است:

\textbf{واقع‌گرایی سیاسی (Realpolitik)} معتقد است که سیاست باید بر اساس واقعیت‌های قدرت، منافع، و امکانات عمل کند، نه آرمان‌های اخلاقی.

\textbf{آرمان‌گرایی سیاسی (Idealism)} معتقد است که سیاست باید در خدمت ارزش‌های والا باشد: عدالت، صلح، آزادی، برابری.

این دو رویکرد هم در سطح ملی و هم در سطح بین‌المللی حضور دارند و هر یک از مکاتب سیاسی را به شکلی متفاوت تحت تأثیر قرار داده‌اند.

\section{تعاریف و ریشه‌ها}

\subsection{واقع‌گرایی سیاسی}

\textbf{ماکیاولی} (۱۴۶۹-۱۵۲۷) اولین نظریه‌پرداز مدرن واقع‌گرایی بود. او در \textit{شهریار} نوشت که حاکم باید بداند چگونه «خوب نباشد» وقتی لازم است. اخلاق خصوصی با اخلاق سیاسی متفاوت است.

\textbf{هابز} (۱۵۸۸-۱۶۷۹) استدلال کرد که در «وضع طبیعی»، زندگی «تنها، فقیرانه، کثیف، وحشیانه، و کوتاه» است. دولت برای فرار از این وضع لازم است.

\textbf{هانس مورگنتا} (۱۹۰۴-۱۹۸۰) واقع‌گرایی کلاسیک را در روابط بین‌الملل تدوین کرد. او معتقد بود سیاست بین‌الملل، مانند همهٔ سیاست، تلاش برای قدرت است.

\begin{tcolorbox}[
    enhanced,
    colback=rightblue!8,
    colframe=rightblue,
    boxrule=1.5pt,
    arc=5pt,
    fonttitle=\bfseries,
    title={\rl{🔵 تعریف: واقع‌گرایی سیاسی}}
]
\textbf{واقع‌گرایی سیاسی} رویکردی است که:
\begin{itemize}[nosep=true]
    \item قدرت را محور سیاست می‌داند
    \item منافع ملی را راهنمای عمل می‌داند
    \item اخلاق را در سیاست ثانویه یا بی‌ربط می‌داند
    \item به طبیعت انسان بدبین است
    \item نظام بین‌المللی را آنارشیک (بدون حاکم مرکزی) می‌داند
\end{itemize}

\textbf{نمایندگان:} ماکیاولی، هابز، مورگنتا، کیسینجر
\end{tcolorbox}

\subsection{آرمان‌گرایی سیاسی}

\begin{tcolorbox}[
    enhanced,
    colback=centergreen!8,
    colframe=centergreen,
    boxrule=1.5pt,
    arc=5pt,
    fonttitle=\bfseries,
    title={\rl{🟢 تعریف: آرمان‌گرایی سیاسی}}
]
\textbf{آرمان‌گرایی سیاسی} رویکردی است که:
\begin{itemize}
    \item ارزش‌های اخلاقی را راهنمای سیاست می‌داند
    \item صلح، عدالت، و همکاری را ممکن می‌داند
    \item به توانایی انسان برای بهبود باور دارد
    \item نهادهای بین‌المللی را مفید می‌داند
    \item حقوق بشر را جهانشمول می‌داند
\end{itemize}

\textbf{نمایندگان:} کانت، ویلسون، جان رالز، یورگن هابرماس
\end{tcolorbox}


\textbf{امانوئل کانت} (۱۷۲۴-۱۸۰۴) در \textit{صلح پایدار} (۱۷۹۵) استدلال کرد که جمهوری‌ها کمتر به جنگ می‌روند و فدراسیون جمهوری‌ها می‌تواند صلح جهانی بیاورد.

\textbf{وودرو ویلسون} (۱۸۵۶-۱۹۲۴)، رئیس‌جمهور آمریکا، پس از جنگ جهانی اول از «جامعهٔ ملل» و «حق تعیین سرنوشت» دفاع کرد.

\textbf{لیبرالیسم بین‌المللی} امروزه معتقد است که دموکراسی، تجارت، و نهادهای بین‌المللی می‌توانند صلح بیاورند.


\section{مکاتب سیاسی و این دو گرایش}

\begin{table}[H]
\centering
\begin{tabular}{|>{\columncolor{darkgray!10}}r|c|c|p{5cm}|}
\hline
\rowcolor{darkgray!30}
\textbf{مکتب} & \textbf{واقع‌گرایی} & \textbf{آرمان‌گرایی} & \textbf{توضیح} \\
\hline
\textbf{مارکسیسم} & متوسط & بالا & آرمان جامعهٔ بی‌طبقه؛ اما واقع‌گرایی طبقاتی \\
\hline
\textbf{لنینیسم} & بالا & متوسط & هدف اخلاقی، اما روش‌های واقع‌گرایانه \\
\hline
\textbf{سوسیال دموکراسی} & متوسط & متوسط & ترکیب آرمان و عمل \\
\hline
\textbf{لیبرالیسم کلاسیک} & متوسط & بالا & حقوق طبیعی؛ اما بازار واقع‌گرا \\
\hline
\textbf{نئولیبرالیسم} & بالا & پایین & بازار واقعیت است؛ عدالت ثانویه \\
\hline
\textbf{محافظه‌کاری} & بالا & پایین & طبیعت انسان ثابت؛ آرمانشهر خطرناک \\
\hline
\textbf{لیبرتارینیسم} & متوسط & بالا & آرمان آزادی؛ اما واقع‌گرایی بازار \\
\hline
\textbf{فاشیسم} & بالا & بالا & قدرت واقعی؛ اما آرمان ملت برتر \\
\hline
\textbf{آنارشیسم} & پایین & بالا & آرمان جامعهٔ بدون دولت \\
\hline
\end{tabular}
\caption{گرایش مکاتب به واقع‌گرایی یا آرمان‌گرایی}
\end{table}

\subsection{تحلیل: وقتی ادعا با عمل نمی‌خواند}

\begin{tcolorbox}[
    enhanced,
    colback=leftred!8,
    colframe=leftred,
    boxrule=1.5pt,
    arc=5pt,
    fonttitle=\bfseries,
    title={\rl{🔴 تحلیل انتقادی: تناقضات}}
]

\textbf{مارکسیسم-لنینیسم:}
\begin{itemize}
    \item \textbf{ادعا:} رهایی بشریت
    \item \textbf{عمل:} گولاگ، سرکوب
    \item \textbf{توجیه:} «واقع‌گرایی انقلابی»
\end{itemize}

\textbf{لیبرالیسم غربی:}
\begin{itemize}
    \item \textbf{ادعا:} حقوق بشر جهانشمول
    \item \textbf{عمل:} حمایت از دیکتاتورها
    \item \textbf{توجیه:} «منافع ملی»
\end{itemize}



\textbf{ناسیونالیسم:}
\begin{itemize}
    \item \textbf{ادعا:} آزادی ملت
    \item \textbf{عمل:} سرکوب اقلیت‌ها
    \item \textbf{توجیه:} «یکپارچگی ملی»
\end{itemize}

\textbf{اسلام‌گرایی:}
\begin{itemize}
    \item \textbf{ادعا:} عدالت الهی
    \item \textbf{عمل:} استبداد دینی
    \item \textbf{توجیه:} «حکم خدا»
\end{itemize}


\vspace{3mm}
\textbf{نتیجه:} همهٔ مکاتب، هنگام رسیدن به قدرت، تمایل به واقع‌گرایی پیدا می‌کنند. آرمان‌ها اغلب توجیه‌گر قدرت می‌شوند، نه محدودکنندهٔ آن.
\end{tcolorbox}

\section{نظم بین‌المللی: واقع‌گرایی یا آرمان‌گرایی؟}

\subsection{معاهدات کلیدی و بازیگران}

\begin{longtable}{|p{3cm}|p{2.5cm}|p{2.5cm}|p{2.5cm}|p{3cm}|}
\hline
\rowcolor{rightblue!30}
\textbf{معاهده/نهاد} & \textbf{پیشنهاددهنده} & \textbf{حامیان اصلی} & \textbf{منتقدان} & \textbf{نتیجه} \\
\hline
\endfirsthead

\textbf{جامعهٔ ملل (۱۹۲۰)} & ویلسون (آرمان‌گرا) & لیبرال‌های بین‌المللی & آمریکا (ایزولاسیونیست)، واقع‌گرایان & شکست \\
\hline
\textbf{سازمان ملل (۱۹۴۵)} & آمریکا، بریتانیا & لیبرال‌ها، چپ میانه & واقع‌گرایان (حاکمیت) & موفقیت نسبی \\
\hline
\textbf{اعلامیهٔ حقوق بشر (۱۹۴۸)} & غرب لیبرال & دموکراسی‌ها & شوروی (رأی ممتنع)، عربستان & تأثیرگذار اما ضعیف \\
\hline
\textbf{برتون وودز (۱۹۴۴)} & آمریکا (کینز/وایت) & غرب & جهان سوم (بعداً) & IMF، بانک جهانی \\
\hline
\textbf{ناتو (۱۹۴۹)} & آمریکا، بریتانیا & غرب & شوروی، چپ & موفق (از دید غرب) \\
\hline
\textbf{پیمان ورشو (۱۹۵۵)} & شوروی & بلوک شرق & غرب & منحل شد \\
\hline
\textbf{معاهدهٔ رم/اتحادیهٔ اروپا (۱۹۵۷)} & فدرالیست‌ها، دموکرات‌های مسیحی & لیبرال‌ها، میانه & ناسیونالیست‌ها، چپ رادیکال & موفق تا برگزیت \\
\hline
\textbf{NPT - منع گسترش هسته‌ای (۱۹۶۸)} & آمریکا، شوروی & اکثر کشورها & هند، پاکستان، اسرائیل & نسبتاً موفق \\
\hline
\textbf{پروتکل کیوتو (۱۹۹۷)} & اتحادیهٔ اروپا، جنبش سبز & لیبرال‌ها، چپ & آمریکا (بوش)، راست & ضعیف \\
\hline
\textbf{دادگاه کیفری بین‌المللی (۲۰۰۲)} & لیبرال‌های بین‌المللی & اروپا، چپ & آمریکا، روسیه، چین & محدود \\
\hline
\textbf{توافق پاریس (۲۰۱۵)} & فرانسه، سازمان ملل & اکثر کشورها & ترامپ، پوپولیست‌های راست & شکننده \\
\hline
\end{longtable}

\subsection{الگوهای همکاری و تضاد}

\begin{center}
\resizebox{\columnwidth}{!}{%
\begin{tikzpicture}[scale=1, every node/.style={font=\small}]

% Title
\node[font=\bfseries\large] at (7,8) {\rl{الگوهای همکاری و تضاد ایدئولوژیک در نظام بین‌الملل}};

% Main actors
\node[draw=rightblue, fill=rightblue!20, rounded corners, minimum width=3cm, minimum height=1cm] (US) at (0,5) {\rl{آمریکا (لیبرالیسم)}};

\node[draw=leftred, fill=leftred!20, rounded corners, minimum width=3cm, minimum height=1cm] (USSR) at (7,5) {\rl{شوروی (سوسیالیسم)}};

\node[draw=goldaccent, fill=goldaccent!20, rounded corners, minimum width=3cm, minimum height=1cm] (EU) at (3.5,2) {\rl{اروپا (سوسیال‌لیبرال)}};

\node[draw=centergreen, fill=centergreen!20, rounded corners, minimum width=3cm, minimum height=1cm] (THIRD) at (10.5,2) {\rl{جهان سوم (متنوع)}};

\node[draw=violet, fill=violet!20, rounded corners, minimum width=3cm, minimum height=1cm] (CHINA) at (14,5) {\rl{چین (سوسیالیسم بازار)}};

% Relations
\draw[<->, line width=2pt, color=leftred] (US) -- (USSR) node[midway, above] {\footnotesize\rl{جنگ سرد}};

\draw[<->, line width=1.5pt, color=centergreen] (US) -- (EU) node[midway, left] {\footnotesize\rl{اتحاد}};

\draw[<->, line width=1.5pt, color=goldaccent, dashed] (USSR) -- (THIRD) node[midway, right] {\footnotesize\rl{حمایت}};

\draw[<->, line width=1.5pt, color=goldaccent, dashed] (US) -- (THIRD) node[midway, below] {\footnotesize\rl{رقابت}};

\draw[<->, line width=2pt, color=violet] (US) -- (CHINA) node[midway, above] {\footnotesize\rl{رقابت جدید}};

\end{tikzpicture}%
}
\end{center}

\subsection{کجا همکاری کارساز بود؟}

\begin{tcolorbox}[
    enhanced,
    colback=centergreen!10,
    colframe=centergreen,
    boxrule=1.5pt,
    arc=5pt,
    fonttitle=\bfseries,
    title={\rl{✅ موفقیت‌های همکاری بین ایدئولوژیک}}
]
\begin{enumerate}
    \item \textbf{سازمان ملل:} با همهٔ محدودیت‌ها، از جنگ جهانی سوم جلوگیری کرد.
    \item \textbf{NPT:} گسترش هسته‌ای را (نسبتاً) کنترل کرد.
    \item \textbf{اتحادیهٔ اروپا:} صلح در قاره‌ای که دو جنگ جهانی داشت.
    \item \textbf{پایان آپارتاید:} فشار بین‌المللی + مذاکره.
    \item \textbf{لایهٔ ازون:} پروتکل مونترال موفق‌ترین توافق زیست‌محیطی.
\end{enumerate}

\textbf{الگو:} همکاری وقتی موفق است که: ۱) منافع مشترک وجود داشته باشد، ۲) هزینهٔ عدم همکاری بالا باشد، ۳) نهادهای نظارتی وجود داشته باشد.
\end{tcolorbox}

\subsection{کجا تضاد به خرابی انجامید؟}

\begin{tcolorbox}[
    enhanced,
    colback=leftred!10,
    colframe=leftred,
    boxrule=1.5pt,
    arc=5pt,
    fonttitle=\bfseries,
    title={\rl{❌ شکست‌های ناشی از تضاد ایدئولوژیک}}
]
\begin{enumerate}
    \item \textbf{جنگ سرد:} تریلیون‌ها دلار هدر رفته، جنگ‌های نیابتی.
    \item \textbf{جنگ ویتنام:} میلیون‌ها کشته به نام «مبارزه با کمونیسم».
    \item \textbf{جنگ‌های خاورمیانه:} تقابل ایدئولوژیک ملی‌گرایانه، دینی، و ابرقدرت‌ها.
    \item \textbf{جنگ عراق ۲۰۰۳:} «صدور دموکراسی» با نتایج فاجعه‌بار.
    \item \textbf{بحران اقلیمی:} تضاد منافع مانع اقدام جدی شده.
\end{enumerate}

\textbf{الگو:} تضاد وقتی به خرابی می‌انجامد که: ۱) ایدئولوژی بر منافع واقعی غلبه کند، ۲) گفتگو ناممکن شود، ۳) «دیگری» شیطانی شود.
\end{tcolorbox}

\section{درس‌ها برای امروز و فردا}

\begin{tcolorbox}[
    enhanced,
    colback=goldaccent!10,
    colframe=goldaccent,
    boxrule=2pt,
    arc=8pt,
    fonttitle=\bfseries\large,
    title={\rl{🎯 درس‌های کلیدی}}
]

\textbf{۱. واقع‌گرایی و آرمان‌گرایی هر دو لازم‌اند.}
واقع‌گرایی بدون اخلاق به جنایت می‌انجامد. آرمان‌گرایی بدون واقع‌بینی به شکست. بهترین سیاست ترکیبی از هر دو است.

\textbf{۲. ایدئولوژی ابزار است، نه هدف.}
وقتی ایدئولوژی از ابزار به هدف تبدیل شود، انسان‌ها فدای آن می‌شوند. همهٔ فجایع قرن بیستم این درس را دارند.

\textbf{۳. نهادهای بین‌المللی مهم‌اند.}
با همهٔ نقص‌ها، سازمان ملل و نهادهای مشابه از جنگ‌های بیشتر جلوگیری کرده‌اند.

\textbf{۴. گفتگو با مخالف ضروری است.}
جنگ سرد وقتی پایان یافت که گفتگو ممکن شد. شیطانی کردن «دیگری» راه به جایی نمی‌برد.

\textbf{۵. چالش‌های جهانی نیازمند همکاری جهانی‌اند.}
تغییر اقلیم، همه‌گیری، هوش مصنوعی—این مسائل در چارچوب ملی حل‌شدنی نیستند.

\textbf{۶. منافع و ارزش‌ها می‌توانند همسو شوند.}
صلح به نفع همه است. تجارت منافع مشترک ایجاد می‌کند. حقوق بشر ثبات می‌آورد. واقع‌گرایی روشن‌بینانه می‌تواند به همان نتایجی برسد که آرمان‌گرایی می‌خواهد.

\textbf{۷. قدرت بدون مشروعیت پایدار نیست.}
امپراتوری‌ها بدون رضایت فرومی‌پاشند. حتی واقع‌گراترین سیاستمداران به توجیه اخلاقی نیاز دارند.

\end{tcolorbox}

\section{نتیجه‌گیری: به سوی واقع‌گرایی اخلاقی}


شاید بهترین رویکرد، آنچه می‌توان «واقع‌گرایی اخلاقی» یا «آرمان‌گرایی واقع‌بینانه» نامید:

\begin{itemize}
    \item پذیرش اینکه قدرت واقعیت دارد
    \item اما قدرت باید در خدمت ارزش‌ها باشد
    \item پذیرش محدودیت‌های انسانی
    \item اما تلاش برای بهبود تدریجی
    \item پذیرش منافع ملی
    \item اما تعریف گستردهٔ منافع شامل صلح و همکاری
\end{itemize}

این رویکرد نه ساده‌لوحی آرمان‌گرایانه است نه بدبینی واقع‌گرایانه. بلکه تلاشی است برای ترکیب بهترین عناصر هر دو.

جهان امروز با چالش‌هایی روبروست که نه واقع‌گرایی ناب می‌تواند حل کند نه آرمان‌گرایی ناب: تغییر اقلیم، نابرابری جهانی، خطرات فناوری، و بحران دموکراسی. حل این مشکلات نیازمند ترکیبی از واقع‌بینی دربارهٔ آنچه ممکن است و پایبندی به ارزش‌هایی است که زندگی را ارزشمند می‌سازند.


\vspace{10mm}
\begin{center}
\rule{0.6\textwidth}{1pt}

\vspace{3mm}
\textit{«سیاست هنر ممکن است. اما ممکن را خود ما تعریف می‌کنیم.»}

\vspace{2mm}
--- اتو فون بیسمارک (با تعدیل)

\vspace{3mm}
\rule{0.6\textwidth}{1pt}
\end{center}
